\section{Extended Worked Example: Propagation of State Debt}

The main paper keeps only a summary-level worked example because the core comparative lesson is conceptual rather than numerical. This section restores a longer version of that reasoning so the propagation view can be inspected more concretely.

\subsection{Scenario}

Consider a multi-tenant LLM serving cluster with continuous batching, mixed prompt lengths, shared prefix reuse, and a fixed memory budget per decode pool. The runtime must decide whether to admit a new request whose prompt is long enough to threaten the current KV working set. On first inspection, admitting the request may look locally beneficial: the cluster keeps utilization high, and the prefill stage can start immediately. However, that decision changes not only instantaneous occupancy, but also the future legality and value of shared prefix state.

In the survey's five-field frame, the state object is the shared KV-cache working set plus its prefix-sharing metadata. The control surface is admission coupled with possible eviction or demotion. The coupling path is that a local admission decision can destroy future reuse, force later restoration, and raise downstream queueing pressure. The evaluation boundary is not just instantaneous throughput; it includes TTFT, tail latency, and goodput over the next burst phase. The remaining gap is that many serving stacks still optimize admission or eviction locally without closing the full loop over restoration and future reuse debt.

\subsection{Propagation Trace}

The propagation argument can be decomposed into five steps:
\begin{enumerate}[leftmargin=*]
\item \textbf{Admission pressure rises.} A long-prompt request arrives while the decode pool is already populated with shorter ongoing sessions. The immediate question is whether to admit now or delay.
\item \textbf{Local memory slack shrinks.} If the request is admitted, the runtime widens the active prefill footprint and may force earlier reclamation of KV pages belonging to lower-priority or seemingly colder sessions.
\item \textbf{Reuse structure is damaged.} Some of the reclaimed pages would have supported prefix reuse or cheaper continuation for later requests from the same conversation family or structured workflow.
\item \textbf{Restoration debt appears later.} Once the next burst or follow-up call arrives, the system must repopulate state that was previously available cheaply. The visible effect is not only additional memory traffic, but also queue amplification because the restoration competes with live decode work.
\item \textbf{The original gain reverses.} The admission that looked locally throughput-positive can reduce goodput or inflate p99 latency after the coupling path fully unfolds.
\end{enumerate}

\begin{quotebox}
The key point is not the exact numbers but the temporal ordering of debt: the runtime can spend future reuse value in order to improve present occupancy, and that trade is invisible unless admission, reclamation, and restoration are evaluated on the same boundary.
\end{quotebox}

\subsection{Why This Example Was Shortened in the Main Paper}

This worked example was compressed in the main paper for two reasons. First, the main paper needs only the transferable systems lesson: access, execution, and evolution form one loop, so local wins must be charged against later state debt. Second, the full trace consumes page budget disproportionately because it requires scenario assumptions, lifecycle detail, and alternate branches that are useful for auditability but not strictly necessary for the article's thesis. The supplement is therefore the right home for the more detailed walkthrough.

\subsection{Implications for Main-Text Compression}

The example also shows how the supplement should support later main-paper cuts. If future compression removes more of the main paper's serving-lifecycle detail, this section can absorb:
\begin{itemize}[leftmargin=*]
\item longer admission-versus-restoration narratives,
\item fuller prefix-reuse legality discussion,
\item expanded disturbance scenarios such as burst after warmup or tenant-switch after partial reclamation,
\item optional figure material illustrating state-debt propagation over time.
\end{itemize}
